\section*{Hoofdstuk \ref{ch:b3:spectral} --- Compacte operatoren
en de spectraalstelling}

\begin{solution}{exo:b3:spectral:1}
(a) $T(B)$ is een begrensde deelverzameling van de
eindigdimensionale $\operatorname{im}T$: relatief \omterm{def:b3:spectral:compact}{compact} volgens
Heine--Borel (\cref{cor:b3:topology:heineborel}, overgezet met
een lineair \omterm{def:b3:topology:continuity}{homeomorfisme} met $\R^n$).
(b) $I$ \omterm{def:b3:spectral:compact}{compact} betekent dat de gesloten eenheidsbal \omterm{def:b3:spectral:compact}{compact} is,
wat volgens de stelling van Riesz (volume van bachelorjaar 2)
precies in eindige dimensie gebeurt. Had een \omterm{def:b3:topology:compact}{compacte} $T$ een
begrensde inverse $T^{-1}$, dan zou $I = T^{-1}T$ \omterm{def:b3:spectral:compact}{compact} zijn
(\cref{prop:b3:spectral:ideal}): onmogelijk in oneindige
dimensie.
\end{solution}

\begin{solution}{exo:b3:spectral:2}
(a) $\norm{Tx}^2 = \sum\abs{\lambda_n}^2\abs{x_n}^2 \leq
\sup\abs{\lambda_n}^2\norm x^2$, met bijna gelijkheid op de $e_n$
die het supremum verwezenlijken.
(b) ($\Leftarrow$) De afknottingen $T_N$ (behoud $n \leq N$,
daarbuiten nul) hebben eindige rang met $\vertiii{T - T_N} =
\sup_{n>N}\abs{\lambda_n} \to 0$: \omterm{def:b3:spectral:compact}{compact} volgens
\cref{prop:b3:spectral:ideal}. ($\Rightarrow$) Is
$\abs{\lambda_{n_k}} \geq \delta > 0$ langs een deelrij, dan is
$\norm{Te_{n_k} - Te_{n_l}}^2 = \abs{\lambda_{n_k}}^2 +
\abs{\lambda_{n_l}}^2 \geq 2\delta^2$: geen enkele convergente
deelrij van $(Te_{n_k})$.
(c) $T^*$ is de diagonaaloperator met $(\bar\lambda_n)$:
\omterm{def:b3:spectral:selfadjoint}{zelftoegevoegd} dan en slechts dan als alle $\lambda_n \in \R$.
Dan is de standaardbasis $(e_n)$ een orthonormale basis van
eigenvectoren, met eigenwaarden $\lambda_n \to 0$: de
spectraalstelling woordelijk.
\end{solution}

\begin{solution}{exo:b3:spectral:3}
Zij $(x_n)$ begrensd. Dan is $(Bx_n)$ begrensd ($\vertiii B <
\infty$); de \omterm{def:b3:topology:compact}{compactheid} van $S$ haalt er $SBx_{n_k} \to y$ uit,
en de \omterm{def:b3:topology:continuity}{continuïteit} van $A$ geeft $ASBx_{n_k} \to Ay$: dus is
$ASB$ \omterm{def:b3:spectral:compact}{compact}. Is $ST = TS = I$ met $T$ \omterm{def:b3:spectral:compact}{compact} en $\dim H =
\infty$, dan zou $I = ST$ \omterm{def:b3:spectral:compact}{compact} zijn, in tegenspraak met
\cref{exo:b3:spectral:1}(b) --- dezelfde uitspraak, van de andere
kant bekeken.
\end{solution}

\begin{solution}{exo:b3:spectral:4}
(a) Met Cauchy--Schwarz in $y$: $\abs{T_kf(x)}^2 \leq
\bigl(\int\abs{k(x,y)}^2\dd y\bigr)\norm f_2^2$; integreer naar
$x$ (Tonelli): $\norm{T_kf}_2 \leq \norm k_{L^2(\square)} \norm
f_2$.

(b) De familie $e_{mn}(x,y) = e_m(x)\overline{e_n(y)}$ is
orthonormaal in $L^2(\intcc01^2)$ (Tonelli splitst de dubbele
integraal). Totaal: staat $h$ loodrecht op alle $e_{mn}$, dan
staat voor elke $m$ de functie $y \mapsto \int
h(x,y)\overline{e_m(x)}\dd x$ (in $L^2$ volgens Cauchy--Schwarz
en Tonelli) loodrecht op elke $\overline{e_n}$ --- en de
toegevoegden $(\overline{e_n})$ vormen een \omterm{def:b3:hilbert:onb}{hilbertbasis} zodra
$(e_n)$ dat doet (toevoeging is een isometrische bijectie van
$L^2$ die orthogonaliteit en totaliteit behoudt) --- dus is ze
$0$ b.o.; dan staat voor b.o.\ $y$ $h(\cdot, y)$ loodrecht op
elke $e_m$: dus $h(\cdot, y) = 0$ b.o., en $h = 0$ (Tonelli).
Bijgevolg is $(e_{mn})$ een \omterm{def:b3:hilbert:onb}{hilbertbasis}; ontwikkel $k = \sum
c_{mn}e_{mn}$. De afknotting $k_N$ (indices $\leq N$) geeft een
$T_{k_N}$ van eindige rang (beeld in $\operatorname{Vect}(e_1,
\dots, e_N)$), en volgens (a) is
\[
\vertiii{T_k - T_{k_N}} \leq \norm{k - k_N}_{L^2} \to 0 :
\]
$T_k$ is een limiet in norm van operatoren van eindige rang, dus
\omterm{def:b3:spectral:compact}{compact} (\cref{prop:b3:spectral:ideal}).
\end{solution}

\begin{solution}{exo:b3:spectral:5}
(a) Op een eigenvector is $\lambda\norm x^2 = \langle x,
Tx\rangle \geq 0$. De formule is
\cref{prop:b3:spectral:sanorm} met alle waarden $\langle x,
Tx\rangle \geq 0$: de absolute waarde is overbodig.
(b) $(x, y) \mapsto \langle x, Ty\rangle$ is een hermitische
positieve (mogelijk ontaarde) sesquilineaire vorm; het
gebruikelijke bewijs van Cauchy--Schwarz (ontwikkel $\langle x +
t\eu^{\iu\theta}y, T(x + t\eu^{\iu\theta}y)\rangle \geq 0$ en
neem de discriminant) gebruikt de definietheid nergens.
\end{solution}

\begin{solution}{exo:b3:spectral:6}
(a) $\norm{Mf}_2 \leq \norm f_2$, en op $f_n = \sqrt n\, \mathbf
1_{\intcc{1 - 1/n}1}$ (eenheidsvectoren) is $\norm{Mf_n}_2 \geq 1
- \frac1n$: dus $\vertiii M = 1$; \omterm{def:b3:spectral:selfadjoint}{zelftoegevoegd} omdat de
vermenigvuldiger reëel is. Eigenwaarden: $xf(x) = \lambda f(x)$
b.o.\ dwingt $f = 0$ b.o.\ af buiten de nulverzameling $\{x =
\lambda\}$: dus $f = 0$ in $L^2$.
(b) Met dezelfde $f_n$: $\norm{Mf_n - f_n}_2 \leq \frac1n \to 0$,
terwijl $f_n \rightharpoonup 0$ (voor vaste $g \in L^2$ is
$\abs{\langle g, f_n\rangle} \leq \norm{g\,\mathbf
1_{\intcc{1-1/n}1}}_2 \to 0$ volgens de gedomineerde
convergentie). Zou $Mf_{n_k} \to h$ in norm, dan is $f_{n_k} \to
h$, wat $h = 0$ afdwingt (zwakke limiet) terwijl $\norm h = 1$:
dus geen enkele convergente deelrij.
(c) In \cref{lem:b3:spectral:existence} gebruikt precies de
extractie ``$Tx_{n_k} \to y$'' de \omterm{def:b3:topology:compact}{compactheid}; voor $M$
concentreren de maximaliserende rijen zich nabij $x = 1$ en
convergeren hun beelden zwak naar $0$, nooit in norm: de
eigenvector aan de top van het numerieke bereik bestaat
eenvoudigweg niet.
\end{solution}

\begin{solution}{exo:b3:spectral:7}
(a) $V = T_k$ met $k(x, y) = \mathbf 1_{y < x} \in
L^2(\intcc01^2)$: \omterm{def:b3:spectral:compact}{compact} (\cref{exo:b3:spectral:4}). Haar
toegevoegde is de kernoperator met kern $\overline{k(y, x)} =
\mathbf 1_{y > x}$: $V^*f(x) = \int_x^1f$; en $V \neq V^*$ (test
op $f = \mathbf 1$).
(b) Is $Vf = \lambda f$ met $\lambda \ne 0$, dan is $Vf$ \omterm{def:b3:topology:continuity}{continu}
op $\intcc01$ (gedomineerde convergentie in $\int_0^x f$), dus
heeft $f = \frac1\lambda Vf$ een \omterm{def:b3:topology:continuity}{continue} vertegenwoordiger; dan
is $Vf$ $\mathcal C^1$ (de hoofdstelling van de
integraalrekening voor \omterm{def:b3:topology:continuity}{continue} integranden), dus is $f$
$\mathcal C^1$, met $\lambda f' = f$ en $f(0) = \frac1\lambda
Vf(0) = 0$: dus $f = C\eu^{x/\lambda}$ met $C = f(0) = 0$.
(c) $V$ is \omterm{def:b3:spectral:compact}{compact} zonder enige eigenwaarde behalve mogelijk $0$
($Vf = 0$ dwingt $f = 0$ b.o.\ af door de integraal te
differentiëren --- dus zelfs $0$ niet): de spectrale machinerie
vereist werkelijk \omterm{def:b3:spectral:selfadjoint}{zelftoegevoegdheid}, niet enkel \omterm{def:b3:topology:compact}{compactheid}.
\end{solution}

\begin{solution}{exo:b3:spectral:8}
Schrijf $x = \sum_ic_ie_i + z$ met $z \in \ker T$, zodat $\langle
x, Tx\rangle = \sum_i\mu_i\abs{c_i}^2$. \emph{Ondergrens}: op de
eenheidssfeer van $V_k = \operatorname{Vect}(e_1, \dots, e_k)$ is
$\langle x, Tx\rangle = \sum_{i\leq k}\mu_i\abs{c_i}^2 \geq
\mu_k$: het maximum over $V$ van het minimum is dus $\geq \mu_k$.
\emph{Bovengrens}: zij $\dim V = k$ en $W_k =
\overline{\operatorname{Vect}}(e_k, e_{k+1}, \dots) + \ker T$,
van codimensie $k - 1$ (haar \omterm{thm:b3:hilbert:decomposition}{orthogonaal complement} is
$V_{k-1}$); dan is $V \cap W_k \neq \{0\}$ (een lineaire
afbeelding $V \to H/W_k \cong V_{k-1}$ van rang $\leq k - 1$
heeft een kern $\neq 0$), en een eenheidsvector $x \in V\cap W_k$
heeft $\langle x, Tx\rangle = \sum_{i \geq k}\mu_i\abs{c_i}^2
\leq \mu_k$: het minimum over $V$ is dus $\leq \mu_k$. Samen: de
eerste formule; de tweede wordt symmetrisch bewezen (test $W =
W_k$; voor willekeurige $W$ van codimensie $k-1$ geeft $W \cap
V_k \neq 0$ een eenheidsvector met $\langle x, Tx\rangle \geq
\mu_k$). Monotonie: $\langle x, Tx\rangle \leq \langle x,
Sx\rangle$ draagt puntsgewijs door de $\max\min$ heen.
\end{solution}

\begin{solution}{exo:b3:spectral:9}
Herschrijf haar als $f - \lambda Tf = g$. Voor $\lambda = 0$ is
$f = g$, altijd uniek \omterm{def:b3:groups:derived}{oplosbaar}. Voor $\lambda \neq 0$ luidt ze
$(T - \frac1\lambda)f = -\frac g\lambda$, en volgens het
\omterm{thm:b3:spectral:fredholm}{alternatief van Fredholm} (\cref{thm:b3:spectral:fredholm}) met de
eigenwaarden $\lambda_n = \bigl((n + \frac12)\pi\bigr)^{-2}$ van
$T$ (\cref{ex:b3:spectral:minxy}) is ze voor elke $g$ uniek
\omterm{def:b3:groups:derived}{oplosbaar} dan en slechts dan als $\frac1\lambda \neq \lambda_n$
voor elke $n$, dat wil zeggen
\[
\lambda \neq \Bigl(n + \tfrac12\Bigr)^2\pi^2
\qquad (n = 0, 1, 2, \dots).
\]
Bij een uitzonderlijke $\lambda = (n+\frac12)^2\pi^2$ bestaan er
oplossingen dan en slechts dan als $g \perp
\sin\bigl((n{+}\frac12)\pi x\bigr)$, en dan zijn ze uniek op het
optellen van veelvouden van die sinus na.
\end{solution}

\begin{solution}{exo:b3:spectral:10}
(a) $Sx = \lambda x$: coördinaten vergelijken geeft $0 = \lambda
x_1$ en $x_n = \lambda x_{n+1}$; is $\lambda \ne 0$, dan is $x_1
= 0$ en inductief $x = 0$; is $\lambda = 0$, dan dwingt $Sx = 0$
$x = 0$ af ($S$ is isometrisch). Geen eigenwaarden. $S^*x =
\lambda x$ luidt $x_{n+1} = \lambda x_n$: $x = x_1(1, \lambda,
\lambda^2, \dots)$, in $\ell^2$ precies wanneer $\abs\lambda <
1$: een volle open schijf eigenwaarden.
(b) $\norm{Se_n - Se_m} = \norm{e_{n+1} - e_{m+1}} = \sqrt2$: het
beeld van de begrensde $(e_n)$ heeft geen enkele cauchydeelrij;
en evenzo $S^*e_{n+1} = e_n$. Geen van beide is \omterm{def:b3:spectral:compact}{compact}.
(c) Voor \omterm{def:b3:topology:compact}{compacte} \omterm{def:b3:spectral:selfadjoint}{zelftoegevoegde operatoren} vangen de
eigenwaarden alles (\cref{thm:b3:spectral:spectral}); laat je een
van beide hypothesen vallen, dan kunnen de eigenwaarden \omterm{def:b3:complete:complete}{volledig}
verdwijnen ($S$, Volterra) of een schijf vullen ($S^*$): het
robuuste object is het spectrum $\{\lambda : T - \lambda I
\text{ niet inverteerbaar}\}$, waarvan de theorie tot een latere
cursus behoort.
\end{solution}

\begin{solution}{exo:b3:spectral:11}
(a) $S$ is de diagonaaloperator met coëfficiënten $\sqrt{\mu_n}
\to 0$: \omterm{def:b3:spectral:compact}{compact} (\cref{exo:b3:spectral:2}(b), overgezet naar de
basis $(e_n)$ aangevuld met $\ker T$, waar $S = 0$),
\omterm{def:b3:spectral:selfadjoint}{zelftoegevoegd} (reële diagonaal), positief ($\langle x, Sx\rangle
= \sum\sqrt{\mu_n}\abs{\langle e_n, x\rangle}^2$), en $S^2 = T$
term voor term.

(b) $R$ commuteert met $T = R^2$; voor een eigenvector $x$ van
$T$ bij eigenwaarde $\mu$ is $T(Rx) = RTx = \mu Rx$, dus is de
eindigdimensionale eigenruimte $E_\mu = \ker(T - \mu)$
$R$-stabiel. Op $E_\mu$ is $R$ symmetrisch en positief met $R^2 =
\mu\,\mathrm{id}$: haar eigenwaarden $\rho$ voldoen aan $\rho^2 =
\mu$ met $\rho \geq 0$, dus zijn ze alle gelijk aan $\sqrt\mu$,
en een diagonaliseerbare operator met één enkele eigenwaarde is
scalair: $R = \sqrt\mu\,\mathrm{id}$ op $E_\mu$. Op $\ker T$:
$\norm {Rx}^2 = \langle x, R^2x\rangle = \langle x, Tx\rangle =
0$. Dus valt $R$ samen met $S$ op $\ker T$ en op elke eigenruimte,
waarvan het gesloten opspansel $H$ is (spectraalstelling): $R =
S$.

(c) $\sqrt G$ werkt als $\frac1{n\pi}$ op $e_n =
\sqrt2\sin(n\pi x)$: het is de kernoperator met
\[
k(x, y) = \sum_{n\geq1}\frac{2\sin(n\pi x)\sin(n\pi
y)}{n\pi},
\]
waarbij de reeks convergeert in $L^2(\intcc01^2)$ (de
coëfficiënten $\frac1{n\pi}$ liggen in $\ell^2$; de kernen van de
partiële sommen geven de benaderingen van eindige rang). Een
elementaire gesloten vorm is niet nodig: de spectrale kant
\emph{is} de operator.
\end{solution}

\begin{solution}{exo:b3:spectral:12}
(a) $T^*T$ is \omterm{def:b3:spectral:compact}{compact} (product van een begrensde en een \omterm{def:b3:spectral:compact}{compacte
operator}, \cref{exo:b3:spectral:3}), \omterm{def:b3:spectral:selfadjoint}{zelftoegevoegd} ($(T^*T)^* =
T^*T$) en positief ($\langle x, T^*Tx\rangle = \norm{Tx}^2$).
Kern: $T^*Tx = 0 \Rightarrow \norm{Tx}^2 = \langle x,
T^*Tx\rangle = 0 \Rightarrow Tx = 0$, en omgekeerd: dus $\ker
T^*T = \ker T$. De spectraalstelling levert de orthonormale
$(e_n)$ met $T^*Te_n = s_n^2e_n$ en $s_n > 0$, die $(\ker
T)^\perp$ opspannen.

(b) $\langle f_m, f_n\rangle = \frac{\langle Te_m,
Te_n\rangle}{s_ms_n} = \frac{\langle e_m,
T^*Te_n\rangle}{s_ms_n} = \frac{s_n^2}{s_ms_n}\delta_{mn} =
\delta_{mn}$. Ontwikkel $x = x_0 + \sum_n\langle e_n, x\rangle
e_n$ met $x_0 \in \ker T$ (\omterm{thm:b3:hilbert:parseval}{Parseval} op het gesloten opspansel
plus de kern); de \omterm{def:b3:topology:continuity}{continue} $T$ toepassen geeft
\[
Tx = \sum_n\langle e_n, x\rangle\,Te_n =
\sum_ns_n\langle e_n, x\rangle\,f_n,
\]
waarbij de reeks convergeert omdat haar partiële sommen cauchy
zijn ($\norm{\sum_{N<n\leq M}s_n\langle e_n, x\rangle f_n}^2 =
\sum s_n^2\abs{\langle e_n, x\rangle}^2$, gedomineerd door
$\sup_{n>N}s_n^2\cdot\norm x^2$, en $s_n \to 0$).

(c) $\norm{Tx}^2 = \sum_ns_n^2\abs{\langle e_n, x\rangle}^2 \leq
(\max s_n)^2\norm x^2$, bereikt op de maximaliserende $e_n$: dus
$\vertiii T = \max s_n$. De ontbinding op rang $N$ afknotten laat
een operator van norm $\sup_{n>N}s_n \to 0$ over: benadering door
eindige rang. De volterra-operator heeft helemaal geen
eigenwaarden (\cref{exo:b3:spectral:7}), maar $V^*V$ wel:
$V^*Vf(x) = \int_x^1\int_0^tf(s)\,\dd s\,\dd t$ is een
symmetrische positieve kernoperator (kern $1 - \max(x,y)$, een
greenkern van het snaartype), waarvan de eigenparen --- berekend
met het randwaardeprobleem $-u'' = \lambda^{-1}u$, $u'(0) = u(1)
= 0$, dat wil zeggen de familie van \cref{exo:b3:spectral:9} ---
de singuliere waarden $s_n = \bigl((n + \frac12)\pi\bigr)^{-1}$
geven. De ontbinding leeft op \emph{twee} orthonormale families
juist omdat $V$ haar eigenmeetkunde wegdraait: geen
eigenvectoren, en toch een volmaakte diagonaalstructuur tussen
twee verschillende bases.
\end{solution}

\begin{solution}{pb:b3:spectral:1}
\textbf{1.} \omterm{def:b3:topology:continuity}{Continuïteit}: $\min$ en $\max$ zijn \omterm{def:b3:topology:continuity}{continu};
symmetrie: $x$ en $y$ verwisselen verandert noch $\min(x,y)$ noch
$1 - \max(x,y)$. Grenzen: $0 \leq g$, en afschattingen van het
type $\max\cdot(1-\max)$ geven $g \leq \frac14$ (met $u = \max$:
$\min \leq u$, dus $g \leq u(1-u) \leq \frac14$). Verder is $g
\in L^2(\square)$: hilbert-schmidt, dus $G$ \omterm{def:b3:spectral:compact}{compact}
(\cref{exo:b3:spectral:4}); en de kern is reëel symmetrisch: $G$
is \omterm{def:b3:spectral:selfadjoint}{zelftoegevoegd}.

\textbf{2.} Splitsen bij $y = x$:
\[
u(x) = (1 - x)\int_0^x y\,f(y)\,\dd y +
x\int_x^1(1 - y)\,f(y)\,\dd y .
\]
Differentieer voor \omterm{def:b3:topology:continuity}{continue} $f$ (productregel en de
hoofdstelling):
\[
u'(x) = -\int_0^xyf + \int_x^1(1-y)f
\qquad\text{(de randtermen vallen weg)},
\]
en $u''(x) = -xf(x) - (1 - x)f(x) = -f(x)$; duidelijk is $u(0) =
u(1) = 0$. Omgekeerd voldoet, als $u \in \mathcal C^2$ in beide
uiteinden verdwijnt, $w = u - G(-u'')$ aan $w'' = 0$ en $w(0) =
w(1) = 0$: $w$ is affien en verdwijnt tweemaal, dus $w = 0$.

\textbf{3.} Zij $Gf = 0$ met $f \in L^2$. Voor $\psi \in \mathcal
C_c^\infty(\intoo01)$ is $\psi = G(-\psi'')$ volgens vraag 2,
dus
\[
\langle f, \psi\rangle = \langle f, G(-\psi'')\rangle
= \langle Gf, -\psi''\rangle = 0
\]
($G$ is \omterm{def:b3:spectral:selfadjoint}{zelftoegevoegd}): volgens het hoofdlemma
(\cref{cor:b3:lp:fundlemma}) is $f = 0$ b.o. Positiviteit: voor
\omterm{def:b3:topology:continuity}{continue} $f$ is, met $u = Gf$,
\[
\langle f, Gf\rangle = \int_0^1 fu = \int_0^1(-u'')u
= \bigl[-u'u\bigr]_0^1 + \int_0^1(u')^2 = \int_0^1(u')^2
\geq 0 ;
\]
voor $f \in L^2$ benader je in $L^2$ door \omterm{def:b3:topology:continuity}{continue} $f_n$: beide
leden gaan naar de limiet ($G$ is begrensd).

\textbf{4.} Is $Gu = \lambda u$ met $\lambda \neq 0$, dan is $Gu$
\omterm{def:b3:topology:continuity}{continu} ($\abs{Gu(x) - Gu(x')} \leq \norm{g(x,\cdot) -
g(x',\cdot)}_2\norm u_2$, en de kern is uniform \omterm{def:b3:topology:continuity}{continu}), dus
heeft $u$ een \omterm{def:b3:topology:continuity}{continue} vertegenwoordiger; de formules van vraag 2
tonen dan dat $Gu \in \mathcal C^2$, dus $u = \frac1\lambda Gu
\in \mathcal C^2$ met $-\lambda u'' = -(Gu)'' = u$ en $u(0) =
u(1) = 0$.

\textbf{5.} $-\lambda u'' = u$ met $u(0) = 0$ geeft $u =
A\sin(x/\sqrt \lambda)$ ($\lambda$ positief: volgens vraag 3 is
$\lambda = \langle u, Gu\rangle/\norm u^2 > 0$ op eigenvectoren).
$u(1) = 0$ dwingt $\frac1{\sqrt\lambda} = n\pi$ af: $\lambda_n =
\frac1{n^2\pi^2}$, met eigenfuncties $\sin(n\pi x)$, genormeerd
$e_n = \sqrt2\sin(n\pi x)$ (want $\int_0^12\sin^2(n\pi x)\dd x =
1$). Controle van de orthogonaliteit: $2\sin(m\pi x)\sin(n\pi x)
= \cos((m-n)\pi x) - \cos((m+n)\pi x)$ integreert voor $m \neq n$
tot $0$.

\textbf{6.} $\ker G = \{0\}$ (vraag 3), dus geeft
\cref{thm:b3:spectral:spectral}(1) dat $H =
\overline{\operatorname{Vect}}(e_n)$: de sinussen vormen een
\omterm{def:b3:hilbert:onb}{hilbertbasis} van $L^2(\intcc01)$. Voor $f(x) = x(1 - x)$:
\[
c_n = \sqrt2\int_0^1x(1-x)\sin(n\pi x)\,\dd x
= \sqrt2\;\frac{2\bigl(1 - (-1)^n\bigr)}{n^3\pi^3}
= \begin{cases}\dfrac{4\sqrt2}{n^3\pi^3} & n \text{ oneven},\\
0 & n \text{ even},\end{cases}
\]
(twee partiële integraties). \omterm{thm:b3:hilbert:parseval}{Parseval}: $\int_0^1x^2(1-x)^2\dd x =
\frac1{30} = \sum_{n \text{ oneven}}\frac{32}{n^6\pi^6}$, dat wil
zeggen $\sum_{n\text{ oneven}}n^{-6} = \frac{\pi^6}{960}$.

\textbf{7.} $\langle e_n, Ge_n\rangle = \lambda_n\norm{e_n}^2 =
\lambda_n$. Voor vaste $x$ zijn de coëfficiënten van $g(x,
\cdot)$: $\langle e_n, g(x,\cdot)\rangle = (Ge_n)(x) =
\lambda_ne_n(x)$, dus $g(x,\cdot) = \sum_n\lambda_ne_n(x)\,e_n$
in $L^2$. De expliciete reeks $\sum_n\lambda_ne_n(x)e_n(y) =
\sum_n\frac{2\sin(n\pi x)\sin(n\pi y)}{n^2\pi^2}$ convergeert
normaal op het vierkant ($\abs{\text{term}} \leq
\frac2{n^2\pi^2}$): haar som is \omterm{def:b3:topology:continuity}{continu}, en ze heeft voor elke
$x$ dezelfde coëfficiënten in $L^2(\dd y)$ als $g(x, \cdot)$: de
twee \omterm{def:b3:topology:continuity}{continue} functies vallen dus in elk punt $(x, y)$ samen.
Stel $y = x$ en integreer (de normale convergentie laat term voor
term integreren toe):
\[
\int_0^1g(x,x)\,\dd x =
\sum_n\lambda_n\int_0^1e_n(x)^2\dd x = \sum_n\lambda_n .
\]

\textbf{8.} $\int_0^1g(x,x)\dd x = \int_0^1x(1 - x)\dd x =
\frac16$, dus $\sum_{n\geq1}\frac1{n^2\pi^2} = \frac16$:
\[
\zeta(2) = \sum_{n\geq1}\frac1{n^2} = \frac{\pi^2}6 .
\]

\textbf{9.} Voor $f = \mathbf 1$ is $c_n =
\sqrt2\int_0^1\sin(n\pi x)\dd x = \sqrt2\,\frac{1 -
(-1)^n}{n\pi}$: dus $c_n = \frac{2\sqrt2}{n\pi}$ voor oneven $n$
en $0$ voor even $n$. \omterm{thm:b3:hilbert:parseval}{Parseval}: $1 = \sum_{n\text{
oneven}}\frac{8}{n^2\pi^2}$, dus $\sum_{n\text{ oneven}}n^{-2} =
\frac{\pi^2}8$, en $\zeta(2) = \frac{\pi^2}8\cdot\frac{1}{1 -
\frac14} = \frac{\pi^2}6$ (de even termen zijn
$\frac14\zeta(2)$). Vergelijking: \omterm{thm:b3:hilbert:parseval}{Parseval} voor één $f$ sommeert
$\abs{\langle e_n, f\rangle}^2$; de spoorformule integreert de
diagonaal van de \emph{kern}, wat neerkomt op \omterm{thm:b3:hilbert:parseval}{Parseval} sommeren
over een hele orthonormale familie tegelijk ---
$\sum_n\langle e_n, Ge_n\rangle$ --- en is daarom blind voor elke
bijzondere keuze van testfunctie.

\textbf{10.} $\abs{c_n\cos(n\pi t)} \leq \abs{c_n}$ met
$\sum\abs{c_n}^2 < \infty$: voor elke $t$ convergeert de reeks in
$L^2$ (orthonormale ontwikkeling,
\cref{thm:b3:hilbert:parseval}(3)); de staartafschatting
$\norm{u(t) - u_N(t)}_2^2 \leq \sum_{n>N}\abs{c_n}^2$ is uniform
in $t$, en elke partiële som is \omterm{def:b3:topology:continuity}{continu} in $t$ (eindig veel
cosinussen): dus is $t \mapsto u(t,\cdot)$ \omterm{def:b3:topology:continuity}{continu} met waarden in
$L^2$. Voor $f = \sum_{n\leq N}c_ne_n$: elke modus $\cos(n\pi
t)\sin(n\pi x)$ voldoet aan $\partial_t^2 = -n^2\pi^2 =
\partial_x^2$ erop toegepast, verdwijnt in $x = 0, 1$, heeft in
$t = 0$ waarde $\sin(n\pi x)$ en tijdsafgeleide $0$: de eindige
som lost alles op. Muzikaal: de beweging van de snaar is een
superpositie van \emph{staande golven} $e_n$, waarvan de
frequenties $n\pi$ de grondtoon en haar boventonen zijn; de
spectraalstelling zegt dat elke beginvorm uniek in deze zuivere
tonen uiteenvalt, met de coëfficiënten $c_n$ als klankkleur. Een
snaar horen is een orthonormale ontwikkeling berekenen.

\textbf{11.} Met $a_n = \abs{\langle u_n, x\rangle}^2$ is
$m_{p+1} = \sum_n\mu_n^{p+1}a_n = \sum_n\bigl(\mu_n^{p/2}
\sqrt{a_n}\bigr)\bigl(\mu_n^{p/2+1}\sqrt{a_n}\bigr) \leq
\sqrt{m_p\,m_{p+2}}$ (Cauchy--Schwarz in $\ell^2$). Bijgevolg
zijn de verhoudingen $m_{p+1}/m_p$ niet-dalend in $p$; omdat
$R(x) = \frac{m_1}{m_0}$, $\frac{\langle Ax, Ax\rangle} {\langle
x, Ax\rangle} = \frac{m_2}{m_1}$ en $R(Ax) = \frac{m_3}{m_2}$,
volgt de keten, met elke term $\leq \mu_1$ omdat $m_{p+1} \leq
\mu_1m_p$ term voor term. Convergentie: is $a_1 > 0$ (waarbij
$a_1$ het totale gewicht van de grootste eigenwaarde is), dan is
\[
\mu_1 \geq R(A^kx) = \frac{m_{2k+1}}{m_{2k}} =
\mu_1\,\frac{a_1 + \sum_{\mu_n<\mu_1}(\mu_n/\mu_1)^{2k+1}
a_n}{a_1 + \sum_{\mu_n<\mu_1}(\mu_n/\mu_1)^{2k}a_n}
\longrightarrow \mu_1,
\]
volgens de gedomineerde convergentie van de sommen (verhoudingen
$< 1$): de machtsmethode convergeert voor elke startvector die
niet loodrecht op de bovenste eigenruimte staat.

\textbf{12.} Puntsgewijs is $\abs{\langle x, (A - B)x\rangle}
\leq \vertiii{A - B}\,\norm x^2$, dus $R_A(x) \leq R_B(x) +
\vertiii{A - B}$ voor elke $x$. Dat invoeren in de
max-minformule van \cref{exo:b3:spectral:8} geeft $\mu_n(A) \leq
\mu_n(B) + \vertiii{A - B}$, en symmetrisch in $A$ en $B$:
$\abs{\mu_n(A) - \mu_n(B)} \leq \vertiii{A - B}$ voor alle $n$
tegelijk.

\textbf{13.} $-w'' = x - x^2$ integreert tot $w = -\frac{x^3}6 +
\frac{x^4}{12} + cx$ (met $w(0) = 0$), en $w(1) = 0$ geeft $c =
\frac1{12}$:
\[
w = \frac{x^4 - 2x^3 + x}{12} = \frac{x(1-x)(1 + x -
x^2)}{12} = Gu .
\]
Verder is $\norm u_2^2 = \int_0^1x^2(1-x)^2 = \frac1{30}$, en met
$\int_0^1x^3(1-x)^3 = B(4,4) = \frac1{140}$:
\[
\langle u, Gu\rangle = \frac1{12}\Bigl(\frac1{30} +
\frac1{140}\Bigr) = \frac{17}{5040},
\qquad
R(u) = \frac{17/5040}{1/30} = \frac{17}{168} .
\]
Dus $\frac1{\pi^2} = \lambda_1 \geq \frac{17}{168}$, dat wil
zeggen $\pi^2 \leq \frac{168}{17} = 9.8824$: $\pi \leq 3.14364 <
3.1437$ (de ware waarde is $\pi^2 = 9.8696$). Eén veelterm, één
integraal, één cijfer.

\textbf{14.} Schrijf $u = x - x^2$, zodat $Gu = \frac{u(1 +
u)}{12}$, en gebruik $\int u^2 = \frac1{30}$, $\int u^3 =
\frac1{140}$ en $\int u^4 = B(5,5) = \frac{4!\,4!}{9!} =
\frac1{630}$:
\[
\norm{Gu}_2^2 = \frac1{144}\int u^2(1+u)^2 =
\frac1{144}\Bigl(\frac1{30} + \frac2{140} +
\frac1{630}\Bigr) = \frac1{144}\cdot\frac{62}{1260} =
\frac{31}{90720} .
\]
Bijgevolg is $\frac{\langle Gu, Gu\rangle}{\langle u, Gu\rangle}
= \frac{31/90720}{17/5040} = \frac{31}{306}$, en volgens vraag 11
is dat nog steeds $\leq \mu_1 = \frac1{\pi^2}$: dus $\pi^2 \leq
\frac{306}{31} = 9.87097$, oftewel $\pi \leq 3.14181 < 3.1419$
--- vier cijfers (en de volgende iteratie zou er ongeveer acht
geven, want de fout krimpt per stap met
$(\lambda_2/\lambda_1)^2 = \frac1{16}$).

\textbf{15.} De familie $(e_m \otimes e_n)(x,y) = e_m(x)e_n(y)$
is een \omterm{def:b3:hilbert:onb}{hilbertbasis} van $L^2(\intcc01^2)$ (orthonormaal volgens
Tonelli; totaal als in \cref{exo:b3:spectral:5}). Volgens vraag 7
is voor vaste $x$ $g(x, \cdot) = \sum_n\lambda_ne_n(x)e_n$, dus
is de coëfficiënt van $g$ op $e_m\otimes e_n$
\[
\langle e_m\otimes e_n,\ g\rangle
= \int_0^1 e_m(x)\,\lambda_n e_n(x)\,\dd x
= \lambda_n\,\delta_{mn} .
\]
\omterm{thm:b3:hilbert:parseval}{Parseval} op het vierkant:
\[
\iint g^2 = \sum_{m,n}\abs{\langle e_m\otimes e_n,
g\rangle}^2 = \sum_n\lambda_n^2 .
\]

\textbf{16.} Wegens de symmetrie van $g$ is
\[
\iint g^2 = 2\iint_{y<x}y^2(1-x)^2
= 2\int_0^1(1-x)^2\,\frac{x^3}3\,\dd x
= \frac23\,B(4, 3) = \frac23\cdot\frac{3!\,2!}{6!}
= \frac23\cdot\frac1{60} = \frac1{90} .
\]

\textbf{17.} Samen: $\sum_n\frac1{n^4\pi^4} = \frac1{90}$, dat
wil zeggen $\zeta(4) = \frac{\pi^4}{90}$. In het algemeen is
$\operatorname{tr}(G^k) = \sum\lambda_n^k =
\frac{\zeta(2k)}{\pi^{2k}}$ gelijk aan een geïtereerde integraal
van producten van de rationale veeltermkern $g$ over de
$k$-kubus: een rationaal getal. Dus is $\zeta(2k) \in \pi^{2k}\Q$
voor elke $k$. De machine bereikt enkel even argumenten omdat de
eigenwaarden erin komen via hun \emph{machten} ---
$\sum\lambda_n^k$ --- en $\lambda_n = \frac1{n^2\pi^2}$: geen
enkele combinatie van sporen brengt $\sum n^{-3}$ voort; de
rekenkundige aard van $\zeta(3)$ (irrationaal volgens Apéry,
transcendentie open) ligt voorbij de spectrale boekhouding.

\textbf{18.} De grootste term van een som van positieve termen is
hoogstens die som: $\lambda_1^2 \leq \sum\lambda_n^2 =
\frac1{90}$, dus $\frac1{\pi^2} \leq \frac1{\sqrt{90}}$ en $\pi
\geq 90^{1/4} = 3.0801\ldots$ Met vraag 14: $3.080 < \pi <
3.1419$, en dat uitsluitend met de rekenkunde van de snaar.
Hogere sporen verscherpen de ondergrens meetkundig: $\lambda_1
\leq (\operatorname{tr}G^{2k})^{1/2k} =
\lambda_1\bigl(1 + \sum_{n\geq2}(\lambda_n/\lambda_1)^{2k}
\bigr)^{1/2k}$, en de parasitaire factor sterft zo snel als
$\bigl(\tfrac14\bigr)^{2k}\cdot\frac1{2k}$ --- hetzelfde
mechanisme van de spectrale kloof als bij de convergentie van de
machtsmethode (vraag 11), nu van de spoorkant bekeken.

\textbf{19.} $\abs{n^2\pi^2 - \nu} \geq \delta > 0$ voor alle $n$
(de rij $n^2\pi^2 \to \infty$ mijdt $\nu$ met een marge), en
$\abs{n^2\pi^2 - \nu} \geq \frac{n^2\pi^2}2$ voor grote $n$.
Convergentie in $L^2$: de coëfficiënten $\frac{c_n}{n^2\pi^2 -
\nu}$ zijn kwadratisch sommeerbaar (gedomineerd door
$\frac{\abs{c_n}}\delta$). Uniforme convergentie: de
supremumnormen van de staarten zijn begrensd door
$\sqrt2\sum_{n>N} \frac{\abs{c_n}}{\abs{n^2\pi^2 - \nu}} \leq
\frac{2\sqrt2}{\pi^2}\bigl(\sum\abs{c_n}^2\bigr)^{1/2}
\bigl(\sum_{n>N}n^{-4}\bigr)^{1/2} \to 0$ (Cauchy--Schwarz).
Verificatie: $Gf + \nu Gu$ heeft als $e_n$-coëfficiënt
\[
\lambda_nc_n + \frac{\nu\lambda_nc_n}{n^2\pi^2 - \nu}
= \frac{c_n}{n^2\pi^2}\Bigl(1 + \frac{\nu}{n^2\pi^2 -
\nu}\Bigr) = \frac{c_n}{n^2\pi^2 - \nu} :
\]
precies de coëfficiënten van $u$, dus $u = Gf + \nu Gu$; en de
uniciteit omdat een verschil $v$ van oplossingen voldoet aan $v =
\nu Gv$, dat wil zeggen $\langle e_n, v\rangle(n^2\pi^2 - \nu) =
0$ voor alle $n$: dus $v = 0$.

\textbf{20.} Als in vraag 19 is de vergelijking $u = Gf + \nu Gu$
gelijkwaardig met de familie coëfficiëntvergelijkingen
$(n^2\pi^2 - \nu)\,\langle e_n, u\rangle = c_n$, $n \geq 1$. Voor
$n = m$ is het linkerlid $0$: de \omterm{def:b3:groups:derived}{oplosbaarheid} dwingt $c_m = 0$
af, en dan is $\langle e_m, u\rangle$ vrij terwijl alle andere
coëfficiënten vastliggen: de oplossingen vormen de rechte $u_0 +
\R e_m$. Resonantie: een aandrijving met een component op de
eigenmodus pompt er onbegrensd energie in --- de schommel,
aangeduwd op haar eigen frequentie.

\textbf{21.} Uit de formule van vraag 19 volgt $\norm{R_\nu
f}_2^2 = \sum\frac{\abs{c_n}^2}{(n^2\pi^2 - \nu)^2} \leq
\frac{\norm f^2}{(\pi^2 - \nu)^2}$ (voor $\nu < \pi^2$ is $\pi^2$
de dichtstbijzijnde eigenwaarde), waarbij de gelijkheid op $f =
e_1$ wordt benaderd: \omterm{def:b3:banach:operator}{operatornorm} $\frac1{\pi^2 - \nu}$.
\omterm{def:b3:topology:compact}{Compactheid}: $R_\nu$ is de limiet in norm van haar afknottingen
van eindige rang (de staartcoëfficiënten $\frac1{n^2\pi^2 - \nu}
\to 0$); de \omterm{def:b3:spectral:selfadjoint}{zelftoegevoegdheid} en de positiviteit lees je van de
diagonaalvorm af (alle coëfficiënten $\frac1{n^2\pi^2 - \nu} >
0$). $R_\nu$ heeft de eigenwaarden $\frac1{n^2\pi^2 - \nu}$: de
analyse van de Delen I tot en met VI begint woordelijk opnieuw.

\textbf{22.} Het woordenboek: \emph{eigenwaarde} $\lambda_n =
\frac1{n^2\pi^2}$ $\leftrightarrow$ de gekwadrateerde frequentie
$n^2\pi^2$ van de $n$-de harmonische; \emph{spoor} $\sum
\lambda_n = \frac16$ $\leftrightarrow$ $\zeta(2) =
\frac{\pi^2}6$; \emph{hilbert-schmidtnorm} $\iint g^2 =
\frac1{90}$ $\leftrightarrow$ $\zeta(4) = \frac{\pi^4}{90}$;
\emph{\omterm{thm:b3:spectral:fredholm}{alternatief van Fredholm}} $\leftrightarrow$ de resonantie
van de aangedreven snaar; \emph{min-max} $\leftrightarrow$
variationele schattingen, tot $\pi < 3.1437$ uit één veelterm.
Achter elk paar hetzelfde object: één \omterm{def:b3:topology:compact}{compacte} \omterm{def:b3:spectral:selfadjoint}{zelftoegevoegde
operator}, één keer gediagonaliseerd, vijf keer uitgebuit.

\textbf{23.} Splitsen naar pariteit en $n = 2m$ substitueren in
het even deel geeft
\[
\zeta(6) = \sum_{n\text{ oneven}}\frac1{n^6} +
\sum_{m\geq1}\frac1{(2m)^6}
= \frac{\pi^6}{960} + \frac{\zeta(6)}{64},
\]
dus $\frac{63}{64}\zeta(6) = \frac{\pi^6}{960}$ en $\zeta(6) =
\frac{64\,\pi^6}{63\cdot960} = \frac{\pi^6}{945}$. De weg via het
spoor zou $\operatorname{tr}G^3 = \sum\lambda_n^3 =
\zeta(6)/\pi^6$ berekenen als $\iint g\,g_2$ met de geïtereerde
kern $g_2(x,y) = \int_0^1g(x,z)g(z,y)\dd z$ --- drie integraties
van stuksgewijze veeltermen; \omterm{thm:b3:hilbert:parseval}{Parseval} op $x(1-x)$ had er maar één
nodig.

\textbf{24.} (a) Diagonaliseer: $v = \sum_na_nu_n$ (plus een
mogelijke component in de kern, waarop $\langle v, Av\rangle$
niets wint terwijl $\norm v^2$ groeit, zodat een maximalisator er
geen heeft). Dan is $\langle v, Av\rangle = \sum\mu_na_n^2 \leq
\mu_1\sum a_n^2$, met gelijkheid dan en slechts dan als $a_n = 0$
zodra $\mu_n < \mu_1$: een maximalisator ligt dus in de
eigenruimte bij $\mu_1$. (b) Voor elke $u$ is
\[
\langle\abs u, A\abs u\rangle - \langle u, Au\rangle
= \iint k(x,y)\,\bigl(\abs{u(x)}\abs{u(y)} -
u(x)u(y)\bigr)\dd x\,\dd y \;\geq\; 0,
\]
want de integrand is puntsgewijs niet-negatief. Hebben $P = \{u >
0\}$ en $N = \{u < 0\}$ allebei positieve \omterm{def:b3:measure:measure}{maat}, dan is de
integrand op $P\times N$ gelijk aan $2k\abs{u(x)}\abs{u(y)} > 0$
op een verzameling van positieve \omterm{def:b3:measure:measure}{maat}: strikte ongelijkheid. Een
eigenfunctie $u$ bij $\mu_1$ maximaliseert het quotiënt van
Rayleigh, en $\abs u$ heeft dezelfde norm, dus zou de striktheid
$R(\abs u) > \mu_1$ opleveren --- onmogelijk; bijgevolg heeft $u$
b.o.\ een constant teken, zeg $u \geq 0$. Dan is $u(x) =
\mu_1^{-1}(Au)(x) = \mu_1^{-1}\int k(x,y)u(y)\dd y > 0$ voor elke
$x \in \intoo01$ (want $k(x,\cdot) > 0$ en $u \neq 0$). (c) Had
de eigenruimte dimensie $\geq 2$, dan zou ze twee
\emph{orthogonale} eigenfuncties $u, v$ bevatten, elk met een
constant teken en volgens (b) inwendig nergens nul; maar dan is
$\abs{\langle u, v\rangle} = \int\abs u\,\abs v > 0$ ---
tegenspraak. Op de snaar: $k = g > 0$ op het open vierkant,
$\mu_1 = \lambda_1 = \frac1{\pi^2}$ is inderdaad enkelvoudig, en
$e_1 = \sqrt2\sin(\pi x)$ is positief op $\intoo01$; en elke
$e_n$ met $n \geq 2$ moet, loodrecht op de positieve $e_1$,
tegen haar in tot nul integreren en dus van teken wisselen ---
zoals haar $n - 1$ inwendige nulpunten $\frac kn$ bevestigen.

\textbf{25.} Omdat $n^2\pi^2 \to \infty$, wordt het minimum $d =
\min_n\abs{n^2\pi^2 - \nu}$ bereikt, in een modus $m$, en $d > 0$
omdat $\nu$ het spectrum mijdt. De diagonaalformule van vraag 19
geeft
\[
\norm{R_\nu f}_2^2 =
\sum_n\frac{\abs{c_n}^2}{(n^2\pi^2 - \nu)^2}
\leq \frac1{d^2}\,\norm f_2^2,
\]
met gelijkheid voor $f = e_m$: $\vertiii{R_\nu} = \frac1d$, het
omgekeerde van de afstand van $\nu$ tot het spectrum --- het
algemene resolventebeginsel, hier in expliciete coördinaten. De
\omterm{def:b3:spectral:selfadjoint}{zelftoegevoegdheid} lees je van de reële diagonaalcoëfficiënten
af; de \omterm{def:b3:topology:compact}{compactheid} volgt als in vraag 21 (de coëfficiënten gaan
naar $0$, dus convergeren de afknottingen van eindige rang in
norm). De prijs van de resonantie: voor $f = e_1$ en $\nu =
(1-\varepsilon)\pi^2$ geeft de formule $u = \frac{c_1}{\pi^2 -
\nu}e_1 = \frac1{\varepsilon\pi^2}e_1$, tegenover het statische
antwoord $Ge_1 = \frac1{\pi^2}e_1$: een versterking met
$\frac1\varepsilon$. Bij $\varepsilon = 10^{-2}$ is het antwoord
$100$ maal het statische --- en het divergeert als $\varepsilon
\to 0$, wat het alternatief van vraag 20 is, gezien van de
begrensde kant: hoe dichter de aandrijffrequentie bij een
eigenfrequentie ligt, hoe minder begrensd de inverse.
\end{solution}
